\section{Preliminaries}
\label{sec:preliminaries}

The coordinate norm $\|\cdot\|_\infty$ on $\mathbb R^S$ is used throughout
because it is simultaneously the norm defining the static ball (7) and the
norm in the gradient envelope (9).  For a distribution $\mathcal D$ on the
finite set $\mathcal X$, $\mathcal D^T$ denotes the law of the independent
training tuple $(x^{(0)},\ldots,x^{(T-1)})$; it is used to keep the conditional
sample averaging explicit.  The law $\mathcal P_{\rm gate}$ and dimension
$d_{\rm path}$ retain the definitions (11)--(12), which are the public feature
objects in the theorem below.  All auxiliary activation envelopes, adjoints,
history-dependent coefficients, and event-split functions are introduced only
inside the appendix arguments where they are derived.
